%% --------------------------------------------------------
\section{Ієрархія методів}
%% --------------------------------------------------------

Post-HF методи утворюють ієрархію за точністю та обчислювальною складністю:

\begin{enumerate}
    \item \textbf{HF} --- базовий рівень, без кореляції
    \item \textbf{MP2} --- $\mathcal{O}(N^5)$ --- найпростіша кореляція
    \item \textbf{MP3, MP4} --- $\mathcal{O}(N^6), \mathcal{O}(N^7)$ --- вищі порядки теорії збурень
    \item \textbf{CCSD} --- $\mathcal{O}(N^6)$ --- надійна динамічна кореляція
    \item \textbf{CCSD(T)} --- $\mathcal{O}(N^7)$ --- "золотий стандарт" квантової хімії
    \item \textbf{Full CI} --- $\mathcal{O}(e^N)$ --- точний розв'язок (у межах базису)
\end{enumerate}


%% --------------------------------------------------------
\subsection{Важливість кореляції для різних властивостей}
%% --------------------------------------------------------

\begin{center}
    \captionof{table}{Вплив електронної кореляції на різні властивості}
    \label{tab:correlation_importance}
    \small
    \begin{tblr}{lll}
        \hline
        \textbf{Властивість}     & \textbf{Важливість кореляції} & \textbf{Рекомендований метод} \\
        \hline
        Абсолютні енергії        & Помірна                       & MP2, DFT                      \\
        Енергії іонізації        & Висока                        & CCSD(T)                       \\
        Електронна спорідненість & Дуже висока                   & CCSD(T) + дифузні             \\
        Енергії збудження        & Висока                        & EOM-CCSD, TD-DFT              \\
        Відстані між станами     & Висока                        & CASPT2, MRCI                  \\
        Дисоціація               & Критична                      & CASSCF, MRCI                  \\
        Слабкі взаємодії         & Критична                      & CCSD(T), MP2-F12              \\
        \hline
    \end{tblr}
\end{center}

\inputcode{code_12.py}


%% --------------------------------------------------------
\subsection{Порівняння методів для He}
%% --------------------------------------------------------

Розглянемо найпростішу багатоелектронну систему --- атом Гелію:

\inputcode{code_13.py}

%% --------------------------------------------------------
\subsection{Порівняння ієрархії методів}
%% --------------------------------------------------------

\inputcode{code_10.py}


Зі зростанням складності методів (\textbf{HF $\to$ MP2 $\to$  CISD $\to$  CCSD $\to$  CCSD(T) $\to$  FCI})
енергія системи послідовно наближається до експериментальної.
Метод \textbf{CCSD(T)} зазвичай відхиляється від експерименту не більш ніж на кілька міліХартрі,
що робить його практичним компромісом між точністю та витратами.



%% --------------------------------------------------------
\subsection{T\textsubscript{1}-діагностика та багатоконфігураційний характер}
%% --------------------------------------------------------

T\textsubscript{1}-діагностика є простим індикатором того, наскільки система є одноконфігураційною.
Вона базується на амплітудах одноелектронних збуджень $t_1$:

\begin{equation}
    T_1 = \frac{\|t_1\|}{\sqrt{N_{\text{elec}}}}
\end{equation}

\inputcode{code_11.py}


Типові значення:
\begin{itemize}
    \item $T_1 < 0.02$ — система добре описується одноконфігураційним наближенням (HF $\Rightarrow$ CCSD точний);
    \item $0.02 < T_1 < 0.045$ — слабка багатоконфігураційність, CCSD залишається надійним;
    \item $T_1 > 0.045$ — сильний багатоконфігураційний характер, CCSD(T) може втрачати точність, слід застосовувати MCSCF або MR-CC.
\end{itemize}

На прикладі \ce{LiH}: при рівноважній відстані $T_1 \approx 0.015$,
що підтверджує одноконфігураційний характер системи.
Під час розтягування зв’язку ($R \to \infty$) $T_1$ зростає, що вказує на
зростання багатоконфігураційності --- саме там CCSD перестає бути повністю точним.
Таким чином, $T_1$ є зручним індикатором придатності методів CC для конкретної системи.
